\documentclass[11pt]{article}
\usepackage[T1]{fontenc}
\usepackage{amsmath,amssymb}
\usepackage[margin=1in]{geometry}
\usepackage[hidelinks]{hyperref}

\title{Origins-DNA v0.2: An Experimental Orbital--Atomic Feature Map}
\author{Adrian Lipa}
\date{15 July 2026}

\newcommand{\DNA}{\mathcal{D}}
\newcommand{\R}{\mathbb{R}}

\begin{document}
\maketitle

\begin{abstract}
This note defines a deterministic feature map for DNA sequences. It is a software-level representation and does not assert a biological mechanism or a decryption of DNA.
\end{abstract}

\section{Definition}
Let \(\DNA=\{A,C,G,T\}\). Assign phases
\begin{equation}
\theta(A)=0,\quad \theta(C)=\frac{\pi}{2},\quad
\theta(T)=\pi,\quad \theta(G)=\frac{3\pi}{2}.
\label{eq:phase-map}
\end{equation}
For each base \(b\in\DNA\), define
\begin{equation}
\Phi(b)=\bigl(\cos\theta(b),\sin\theta(b),r(b),h(b)\bigr)\in\R^4,
\label{eq:feature-map}
\end{equation}
where \(r=+1\) for purines and \(-1\) for pyrimidines, while \(h=-1\) for A/T and \(+1\) for C/G.

\section{Declared symmetry}
For the Watson--Crick complement operator \(c\), the map satisfies
\begin{equation}
\Phi_{1:2}(c(b))=-\Phi_{1:2}(b),\qquad
r(c(b))=-r(b),\qquad h(c(b))=h(b).
\label{eq:complement-symmetry}
\end{equation}
Equation~\eqref{eq:complement-symmetry} is an imposed representation invariant, not an empirical biological result.

\section{Local consolidation}
For a window \(W=(b_1,\ldots,b_n)\), define mean phase components and coherence
\begin{equation}
\bar{x}=\frac{1}{n}\sum_{j=1}^n\cos\theta(b_j),\quad
\bar{y}=\frac{1}{n}\sum_{j=1}^n\sin\theta(b_j),\quad
R_W=\sqrt{\bar{x}^2+\bar{y}^2}.
\label{eq:coherence}
\end{equation}
The implementation also records mean ring class, mean bond class and signed wrapped phase flow.

\section{Validation boundary}
The representation must be compared against one-hot, 2-bit and k-mer baselines under held-out evaluation, label permutation, sequence shuffling and component ablation. Software tests establish implementation correctness only.

\end{document}
